%!TeX root=../main.tex

\chapter{Introdução}
\label{cap:introducao}

    Shogi é um jogo de estratégia japonês semelhante ao xadrez, mas que apresenta 
    mecânicas adicionais que aumentam a dificuldade de seu aprendizado. Entre essas mecânicas, 
    destacam-se a possibilidade de reutilizar peças capturadas e a promoção de peças para versões 
    com movimentos diferentes. 

    Ao longo dos anos, plataformas como o Lishogi surgiram com o intuito de facilitar o acesso ao aprendizado, 
    oferecendo recursos como partidas online ou resolução de problemas como formas de treino. Essas plataformas, 
    no entanto, apresentam limitações que se tornam uma barreira ao aprendizado de iniciantes. O Lishogi, 
    por exemplo, disponibiliza exercícios simples para quem está aprendendo, mas não fornece informações sobre os 
    motivos pelos quais uma resposta está incorreta. O mesmo vale para as análises de partidas, nas quais as jogadas 
    recebem uma pontuação de acordo com sua qualidade, sem explicar por que determinada decisão seria melhor.

    Para um iniciante que deseja aprender a jogar sozinho, a falta de orientação pode dificultar a compreensão dos conceitos envolvidos e assim limitar o aproveitamento dos exercícios realizados. Nesse contexto, o uso de modelos de linguagem pode contribuir para suprir tal necessidade ao fornecer ao jogador explicações moldadas ao seu nível de conhecimento, desempenhando um papel semelhante ao de um tutor. 

    Diante desse cenário, este trabalho propõe o desenvolvimento do Shogi-sensei, uma plataforma web voltada ao ensino
    introdutório de shogi. O projeto fornece exercícios divididos em módulos que abordam os conceitos fundamentais
    do jogo: reconhecimento das peças, movimentação, promoção, reintrodução de peças capturadas e situações de 
    xeque-mate.
    Ao final de cada lista, o usuário recebe os resultados, incluindo as questões respondidas incorretamente e uma análise
    gerada por um modelo de linguagem (LLM) que explica os motivos do erro. A seleção de exercícios de cada lista é adaptada
    ao desempenho do jogador, aumentando a proporção de exercícios dos módulos nos quais ele apresenta maior dificuldade.
    Para avançar para o próximo módulo, é necessário atingir uma média mínima nas listas mais recentes do módulo atual.
 
    A avaliação do sistema foi realizada durante 2 semanas com X usuários. Durante esse período, 
    foram registradas métricas como taxa de acerto, tempo de resposta e distribuição de exercícios por módulo para
    acompanhar o desempenho dos usuários ao longo do experimento.

    O restante deste trabalho está organizado da seguinte forma. O capítulo 2 apresenta os conceitos e tecnologias
    utilizados para desenvolver este trabalho. O capítulo 3 apresenta trabalhos relacionados. O capítulo 4 descreve a implementação
    do sistema. O capítulo 5 contém a metodologia utilizada para avaliar a plataforma desenvolvida. O capítulo 6 
    apresenta e discute os resultados obtidos no experimento descrito no capítulo anterior. Por fim, o capítulo 7
    apresenta as conclusões e possíveis implementações futuras.